% The Prose Barrier --- arXiv Preprint (Expanded)
% Extended from ACL 4-page to arXiv ~10-page format
% arXiv: cs.CL / cs.AI

\documentclass[11pt]{article}
\usepackage[preprint]{acl}

\usepackage{times}
\usepackage{latexsym}
\usepackage[T1]{fontenc}
\usepackage[utf8]{inputenc}
\usepackage{microtype}
\usepackage{inconsolata}

\usepackage{booktabs}
\usepackage{amsmath}
\usepackage{graphicx}
\usepackage{hyperref}

\title{The Prose Barrier: Structural Limits of Agent Self-Verification \\
  and a Five-Layer Configuration Integrity Architecture}

\author{Yuhao Lin \\
  Fujian Agriculture and Forestry University \\
  \texttt{3235706088@stu.fafu.edu.cn} \\\\
  \textbf{Independent undergraduate work. No research supervisor. No funding.}}

\begin{document}
\maketitle

\begin{abstract}
AI coding agents are configured through natural-language rules---system prompts,
markdown files, hook registries. Over long conversations, these rules degrade.
The root cause, I find, is structural: an LLM's generation and self-evaluation
share the same decoder $P(\text{token} \mid \text{context}; \theta)$, so the
model cannot independently verify rule compliance---it can only \textit{generate
a claim} that it did. I call this the Prose Barrier.

I present a five-layer architecture organized by position relative to the
Barrier: L0 psychological safety (uncertainty admission), L1 mechanical gates
(filesystem checks bypassing NL entirely), L2 neural gates (logprob probes
detecting constraint penetration without content interpretation), L3 causal
encoding (syllogistic format altering internal representation depth), and L4
drift prediction (trend-based early warning). Deployed over 34 agent sessions
with 16 experiments across 3 architectures (DeepSeek V4 Pro, Qwen3-8B, GLM-4-9B),
including two pre-registered experiments and one DPO training pipeline.

Three core patterns emerge. First, mechanical gates (L1) reduce rule violations
from 55.9\% to 0.7\%, operating entirely outside the Prose Barrier---no natural
language, no model judgment, just filesystem checks. Second, syllogistic rule
format deepens constraint internalization at the logprob level ($d_z{=}+0.578$,
$BF_{10}{=}2.8{\times}10^5$, objective API-read DV) without changing behavioral
output ($\text{IMP}{\approx}\text{SYL}$ across architectures)---the L2/L3
dissociation. Third, the relationship between output constraint severity and
format effect magnitude follows a non-monotonic pattern ($C_1 > C_0 > C_3 > C_2$),
revealing three processing regimes (optimization, suppression, rebound) rather
than a single slope.

Additional experiments include: a 2$\times$2 factorial (Format $\times$ GateGuard,
$n{=}240$) whose pre-registered hypothesis was not confirmed---format effects on
reasoning depth were independent of gate state; cross-model scanner calibration
revealing that mechanizability classification is architecture-dependent; persona
decorrelation experiments (E1, E1b) testing whether named-persona expert review
provides genuine independence or costume-diversity; and a DPO training pipeline
(Qwen2.5-1.5B, QLoRA on RTX 3060) showing that causal-reasoning preferences can
be partially internalized into model weights (CIS metric, per-token NLL reduction
of 5--10\%).

All experiments by a single undergraduate student on one laptop (Dell G15,
RTX 3060). No supervisor. No funding. Code and data at
\url{https://github.com/YuhaoLin2005/hermes-workspace}.
Independent validation CLI at \url{https://github.com/YuhaoLin2005/paper-validator}.
\end{abstract}

\section{Introduction}

A typical Claude Code agent loads 5--10 markdown files defining its behavioral
rules: how to review code, when to run tests, what constitutes a deliverable.
These rules sit in system prompts and hook registries. The agent reads them at
session start. Then it works---sometimes for hours, across dozens of tasks,
generating thousands of tokens.

Whether the agent actually follows its rules is an empirical question.

I ran 34 sessions with a rule-configured agent. Retrospective coding of these
sessions showed rule violations in 55.9\% of sessions (19 of 34) before
mechanical enforcement was wired in. After adding filesystem-level checks
(mtime comparison, regex matching, exit code verification) and running a
150-task controlled experiment across 6 sessions, the violation rate dropped to
0.7\% (1 of 150 tasks). These behavioral numbers come from a scoring protocol
later found to have poor inter-rater reliability ($\kappa=-0.14$; see
Section~\ref{sec:limitations}), so they should be treated as directional
indicators, not precise estimates.

The scoring protocol's failure was itself informative: the same structural
constraint that prevents the agent from self-verifying (Section~\ref{sec:barrier})
also prevented the single author from blind-scoring reliably---a mirror of the
Prose Barrier at the meta-level, where independent verification requires a second
person rather than a second prompt.

The core question is \textit{why}. A language model that performs at
professional-exam level can fail to reliably verify that it wrote a file to disk.
The answer, I came to believe, is about architecture, not capability.

When an LLM agent evaluates itself---``Did I follow rule $R$ on this task?''---
its answer comes from the same decoder that may have already failed to follow
$R$. There is no independent channel. The act of self-checking and the act of
generating share the same token distribution
$P(\text{token} \mid \text{context}; \theta)$. This is the
\textbf{Prose Barrier}: a structural constraint that makes unassisted agent
self-verification unreliable, not because the model is bad at it, but because
the architecture makes genuine independent verification impossible.

I advance three claims:

\begin{enumerate}
\item \textbf{The Prose Barrier is a structural constraint, not an engineering
  bug.} It follows from the autoregressive nature of LLMs and means that
  verification channels must be architecturally independent of generation
  channels.

\item \textbf{A five-layer architecture systematically addresses this
  constraint.} The layers are organized by their position relative to the
  Prose Barrier---pre-processing (L0), operating outside it (L1), probing inside
  it without interpreting content (L2), altering internal routing through it
  (L3), and observing it from outside across time (L4).

\item \textbf{Format effects on internal representations are measurable at the
  logprob level, but behavioral claims require independent validation.}
  Logprob-level measurement ($d_z=+0.578$, $BF_{10}=282$k) provides objective
  evidence that syllogistic rule format alters internal constraint representations.
\end{enumerate}

\section{The Prose Barrier}
\label{sec:barrier}

\subsection{Definition}

Consider an autoregressive language model $M$ with parameters $\theta$. At
decoding step $n$, the model samples:
\[
t_n \sim P(\cdot \mid t_{1}, \ldots, t_{n-1}; \theta)
\]
Every token the model produces---whether ``thinking about what to do,''
``evaluating what it just did,'' or ``taking action''---is drawn from this same
distribution. There is no separate ``auditor'' module. When the agent is asked
``did you follow rule $R$,'' its answer is sampled from the same decoder whose
state may already encode a failure to follow $R$.

\textbf{Consequence.} The Barrier cannot be circumvented by prompting. Telling
the agent ``be more careful about self-checking'' generates more self-checking
tokens---but those tokens are still drawn from the same distribution.

\textbf{The Prose Barrier of Measurement.} A corollary emerged during
experimentation. To measure format effects via logprobs, I had to constrain the
agent's output (``Output only A or B'')---otherwise it produces paragraphs of
free-form reasoning where the logprob DV cannot be extracted. But this output
constraint tells the model to suppress the reasoning channel that syllogistic
format was designed to enhance. The measurement instrument changes the
phenomenon being measured (Section~\ref{sec:gradient}).

\subsection{The Five-Layer Response}

If the Barrier cannot be removed, it can be worked around:

\begin{description}
\item[L0 --- Psychological Safety.] Pre-processes generation by reframing
  uncertainty. ``Accuracy $>$ completeness. `I don't know' is a valid output.''
  Tested on 40 binary-choice probes: accuracy preserved ($+0.01$),
  non-ceiling uncertainty admission improved ($r=+0.949$, $n=5$).
  \textbf{Status:} preliminary signal only.

\item[L1 --- Mechanical Gate.] Operates entirely outside the Barrier. Instead
  of asking the agent ``did you read back the file?'' (NL, through the Barrier),
  a shell script compares \texttt{mtime} timestamps (filesystem, no NL).
  Deployed over 34 sessions: violation rate 55.9\% $\to$ 0.7\% (exploratory).
  19 automated tests (mtime verification, SHA256 chain integrity, stale output
  detection, exit-code enforcement)---all 19 pass.
  \textbf{Status:} deployed; single-machine only.

\item[L2 --- Neural Gate.] Probes inside the Barrier but avoids content
  interpretation. Measures token logprobs at decision points---if syllogistic
  format produces higher probability for rule-consistent choice tokens vs
  imperative, format penetration is detected without interpreting a single word.
  DV read directly from API JSON (\texttt{logprobs=True}).
  \textbf{Status:} 40-probe validation complete ($d_z=+0.578$, $BF_{10}=282$k);
  single model (DeepSeek V4 Pro).

\item[L3 --- Causal Encoding.] Alters internal processing through the Barrier.
  Reformulating the same rule from imperative (``verify your output'') to
  syllogistic form (``Write operations can deviate from intent. Only a read-back
  confirms correctness. Therefore, verify after writing.'') changes attention
  routing via causal chain structure.
  \textbf{Status:} multiple experiments; causal mechanism relies on Pender (2026).

\item[L4 --- Drift Prediction.] Observes L1--L3 state across time, entirely
  outside the Barrier. 12 mechanical features tracked across sessions.
  \textbf{Status:} deployed; predictive validation pending.
\end{description}

\begin{table}[t]
\centering
\begin{tabular}{lccc}
\toprule
\textbf{Layer} & \textbf{Position} & \textbf{Mechanism} & \textbf{Evidence} \\
\midrule
L0 Safety  & Pre-barrier & Reframe uncertainty   & Accuracy $+0.01$, $r{=}+0.949$ \\
L1 Gate     & Outside     & Filesystem checks     & 55.9\%$\to$0.7\%, 19/19 tests \\
L2 Gate     & Inside      & Logprob probes        & $d_z{=}+0.578$, $BF{=}2.8{\times}10^5$ \\
L3 Encoding & Through     & Syllogistic format    & $d_z$ 0.09--0.60 (non-monotonic) \\
L4 Predict  & Outside     & Drift features        & 12 features deployed \\
\bottomrule
\end{tabular}
\caption{Five-layer architecture by position relative to the Prose Barrier.}
\label{tab:architecture}
\end{table}

\section{Key Experiments}

\subsection{Logprob V3: Does Format Change Internal Representation?}
\label{sec:v3}

An initial attempt ($n=8$ probes) produced $d_z = -0.148$. A DEV.to reader
identified the issue: 4 of 8 probes had comparison tokens absent from DeepSeek's
top-20 logprobs, so the script substituted a $-10.0$ sentinel value---injecting
noise. I built a mechanical pre-screening pipeline (\texttt{probe\_validator.py})
that verifies both option tokens are present before the main experiment runs.

\textbf{Design.} Within-probe, 3-condition (NO RULES / imperative / syllogistic).
40 probes (10 per category: action, epistemic, structural, meta). DV: $\log P(A)
- \log P(B)$ at first decision token, read directly from DeepSeek V4 Pro API
($T=0.2$).

\textbf{Results.} $n=40$. Mean IMP: $-2.95$ ($SD=14.67$). Mean SYL: $+4.48$
($SD=13.45$). $\bar{x}_{\Delta} = +7.43$, $SE=2.03$, $t(39)=3.65$, $p=0.0008$.
Cohen's $d_z = +0.578$ (bootstrap 95\% CI: $[+3.39, +11.17]$). $BF_{10} =
282{,}399$ (JZS prior, $r=0.707$). 32/40 probes (80\%) favoring syllogistic.

\begin{table}[t]
\centering
\begin{tabular}{lrrr}
\toprule
\textbf{Category} & \textbf{IMP} & \textbf{SYL} & $\mathbf{d_z}$ \\
\midrule
Action      & $-3.66$ & $+1.85$ & $+0.50$ \\
Epistemic   & $-4.23$ & $+3.92$ & $+0.76$ \\
Structural  & $-2.37$ & $+5.56$ & $+0.67$ \\
Meta        & $-1.54$ & $+3.21$ & $+0.38$ \\
\bottomrule
\end{tabular}
\caption{Logprob V3 per-category. Category $\times$ Format: $F(3,36)=0.26$, n.s.}
\label{tab:v3}
\end{table}

\subsection{Where Format Helps: L1-Visibility Analysis}
\label{sec:l1vis}

A reviewer predicted format effects should concentrate where mechanical gates
cannot reach. I tested the opposite: format amplifies where L1 provides
structural anchors.

\textbf{Classification.} Each probe classified as L1-visible (violation produces
deterministic mechanical signal) or L1-invisible (violation concerns judgment
quality). Three-test criterion: SIGNAL, ACTION, CERTAINTY. 22/40 L1-visible;
18/40 L1-invisible.

\textbf{Results.} L1-visible: $d_z=+0.71$, 95\% CI $[+0.35, +1.08]$.
L1-invisible: $d_z=+0.40$, 95\% CI $[+0.02, +0.78]$. Between-group $\Delta
d=+0.31$, $t(38)=1.03$, $p=0.31$ (n.s.). Direction consistent with synergy
(format amplifies existing structural anchors), not substitution.

\subsection{Constraint Gradient: A Non-Monotonic Pattern}
\label{sec:gradient}

\textbf{Design.} 12 probes $\times$ 4 constraint levels $\times$ 2 formats = 96
API calls (DeepSeek V4 Pro, $T=0.2$):

\begin{itemize}
\item \textbf{C0}: No constraint (baseline)
\item \textbf{C1}: ``Output only the letter A or B'' (light)
\item \textbf{C2}: ``Output only one letter, no other text, no explanation'' (moderate)
\item \textbf{C3}: ``Do not output any characters except A or B.'' (heavy)
\end{itemize}

\textbf{Results.} The monotonic prediction is wrong:

\begin{table}[t]
\centering
\begin{tabular}{lcccc}
\toprule
\textbf{Level} & $\mathbf{n}$ & $\bar{x}$ & $\mathbf{s}$ & $\mathbf{d_z}$ \\
\midrule
C0 (none)     & 11 & $+4.91$  & 15.58 & $\mathbf{0.315}$ \\
C1 (light)    & 12 & $+7.10$  & 11.92 & $\mathbf{0.596}$ \\
C2 (moderate) & 12 & $+1.13$  & 12.40 & $\mathbf{0.091}$ \\
C3 (heavy)    & 12 & $+2.21$  & 7.42  & $\mathbf{0.297}$ \\
\bottomrule
\end{tabular}
\caption{Constraint gradient. Three regimes: optimize, suppress, rebound.}
\label{tab:gradient}
\end{table}

Three processing regimes: \textbf{Optimization} ($C_0 \to C_1$): $d_z$ rises
0.315$\to$0.596---light constraint improves measurement. \textbf{Suppression}
($C_1 \to C_2$): $d_z$ crashes to 0.091---``no explanation'' kills the reasoning
channel. \textbf{Rebound} ($C_2 \to C_3$): $d_z$ recovers to 0.297 with
tightest variance ($s=7.42$)---system prompt becomes sole differentiator under
extreme compression.

\subsection{Behavioral Replication Across Architectures}
\label{sec:crossmodel}

Same 12 probes with GateGuard-OFF vs NO RULES on 3 models. Compliance (NO RULES
$\to$ IMP $\to$ SYL): DeepSeek V4 Pro (MoE): 0.48 $\to$ 0.86 $\to$ 0.83;
Qwen3-8B (Dense): 0.79 $\to$ 1.00 $\to$ 0.98; GLM-4-9B: 0.83 $\to$ 1.00 $\to$
1.00 (exploratory; $\kappa=-0.14$). IMP $\approx$ SYL on all models ($|\Delta|
\leq 0.03$)---format changes representations, not behavior.

\subsection{P1-1: Residual Violation Clustering}
\label{sec:p1-1}

\textbf{Design.} 5 task types $\times$ 40 trials = 200 API calls. Deterministic
regex scoring. \textbf{Results (INCONCLUSIVE).} T1 (Format-tag): 100\%, T2
(Section-header): 100\%, T3 (Checklist): 0\% (all 40 violations---prose
checklists instead of markdown \texttt{- [ ]}), T4 (Reasoning): 35\%, T5
(Uncertainty): 42.5\%. T3 exposed a mechanizability classification error---the
scanner labeled it L1, but \texttt{- [ ]} syntax was too narrow.

\subsection{P1-2: Format $\times$ GateGuard Factorial (Pre-Registered)}
\label{sec:p1-2}

\textbf{Design.} 2$\times$2 factorial: Format (code vs prose) $\times$ GateGuard
(on vs off) = 4 conditions $\times$ 2 task types $\times$ 30 trials = 240 API
calls. SHA256 pre-registered hypothesis: format effect on reasoning is LARGER
under GateGuard-OFF.

\textbf{Results: NOT\_CONFIRMED.} Format effect on reasoning: Gate ON
$d=-0.277$, Gate OFF $d=-0.250$ (nearly identical). Format effect on mechanical:
Gate ON $d=2.960$ (massive), Gate OFF $d=0.000$ (none). \textbf{The gate
entirely mediates format effects.}

\subsection{P1-1 Cross-Model: Scanner Calibration}
\label{sec:p1-1-cross}

\textbf{Design.} 5 tasks $\times$ 20 trials $\times$ 2 models = 200 calls.
\textbf{Results.} DS Flash: 100\%/95\%/100\%/95\%/25\% (hollow compliance).
Qwen3.6-35B: 40\%/70\%/10\%/40\%/0\%. Scanner alignment: 3/5 tasks agree.
\textbf{Gateability = Rule Structure $\times$ Model Capacity}---2D space.

\subsection{E1: Named-Persona Decorrelation}
\label{sec:e1}

\textbf{Design.} 3 personas (Carmack, Torvalds, Knuth) $\times$ 2 models (DS V4
Pro, Kimi K2.7) $\times$ 5 snippets = 30 calls. Blind LLM-as-judge scoring.

\textbf{Results.} Cross-model same-persona agreement: \textbf{1.000} ($n=15$).
Within-model diff-persona: \textbf{0.867} ($n=30$). Persona drives judgment---
genuine independence, not costume diversity. \textbf{Caveat:} ceiling effect
(4/5 unanimous). \textbf{E1b Replication:} 14 discriminating snippets $\times$ 3
personas $\times$ 2 models $+$ control = 112 calls. Fleiss' $\kappa$ pending.

\subsection{DPO Training: Behavioral Internalization}
\label{sec:dpo}

\textbf{Motivation.} Can rule-following preferences be internalized into weights?

\textbf{Pipeline.} Qwen2.5-1.5B-Instruct, QLoRA ($r{=}64$, $\alpha{=}128$) on
RTX 3060 6GB. DPO: $\beta{=}0.1$, 1 epoch. 200+ preference pairs across 4
domains. \textbf{CIS (Causal Internalization Score)} = OOD $\times$ 0.4 + discrim
$\times$ 0.3 + behavioral $\times$ 0.3 (0--100). PASS: $\geq$20. Per-token NLL
analysis: DPO reduces NLL on causal responses $\sim$5--10\% (Cohen's $d$
computed per domain).

\textbf{Interpretation.} DPO amplifies existing preferences. Modest effects
consistent with Prose Barrier: weight-level internalization faces the same
structural constraint as prompt-level rule-following.

\section{The Dual-Pool Expert Review System}
\label{sec:dual-pool}

\subsection{Design}

\textbf{Fixed Pool.} 33 documented experts across 6 domains, each anchored to
verifiable principles outside the model's generation loop: Torvalds ``eliminate
special cases'' (TED 2016, high), Thompson ``trust boundary'' (1984 Turing
Award, high), Carmack ``measure first'' (.plan files, high), Norman
``affordance'' (Design of Everyday Things, high). Attribution confidence:
high/moderate/low. Rule: no persona below moderate.

\textbf{Random Pool.} Per-session web search for unfamiliar experts. Breaks
fixed pool echo chamber.

\textbf{Routing Table.} YAML task-to-expert mapping. Triggers match tasks to
expert combinations without manual selection.

\subsection{The Core Limitation}

E1 (Section~\ref{sec:e1}) directly tests the central claim: named-persona
review provides genuine independence, not costume diversity. Initial results
promising but ceiling-limited. System explicitly acknowledges: 33 personas on
same model = same inference engine in 33 costumes.

\section{Discussion}

\subsection{What This Work Is and Is Not}

Conceptual contribution: reframing agent reliability from ``make them better''
to ``what makes self-verification structurally impossible.'' Empirical
contribution: demonstrating format effects measurable at logprob level with
objective, API-read DV---bypassing the behavioral scoring problem.

\subsection{Context-Fragility of Format Effects}

Multi-scene P1 experiment ($n=12$): $d_z$ collapsed from 0.58 (single-scene) to
0.19 (multi-scene), 67\% reduction. Bootstrap 95\% CI crosses zero. Per-probe
V3-P1 correlation: $r=-0.65$. Control experiments: meta-instruction alone
drives $\sim$80\% of collapse. Three-parameter model: causal chain complexity
$\times$ cognitive depth $\times$ model capacity.

\subsection{Related Work}

\textbf{LLM Self-Verification Limits.} \citet{kambhampati2024} established LLMs
require external verification modules (LLM-Modulo framework)---L1 instantiates
this for configuration management. \citet{bai2022constitutional} proposed
constitutional AI self-critique; Prose Barrier argues this is structurally
constrained by the shared decoder. \citet{startari2025tloc} independently
proved, via formal theorem (TLOC), that transformer architectures cannot
structurally verify internal rule compliance---the same constraint this work
demonstrates empirically through logprob measurement and behavioral scoring.

\textbf{Runtime Verification Systems.} Multiple systems address agent compliance
through external verification. \citet{sharma2026rvf} built a 4-layer binary
bitmask verification stack (RVF) replacing LLM self-review with deterministic
predicates, finding 660 silent bugs across 261 services---bugs LLM self-review
structurally missed. \citet{kang2026policyguard} inserted a sub-agent verifier
(PolicyGuard) in the tool-calling loop with dialogue-grounded policy reasoning,
achieving +6--12pp Pass\textsuperscript{4} across 3 vendors. Two VIGIL systems
address complementary threats: \citet{lin2026vigil} targets tool stream injection
via speculative verify-before-commit (SIREN benchmark, 959 cases), while
\citet{li2026vigil} enforces behavioral specifications via SMT-based runtime
monitoring ($>$95\% recall, $<$10\% FPR). \citet{zhang2026sentinelshield} provides a
6-shield defense-in-depth middleware with persona drift detection. These systems
share L1's design intuition---external enforcement outside the generation loop---
but none combine logprob-level constraint measurement, non-monotonic constraint
gradient analysis, and DPO-based behavioral internalization within a single
barrier-centered architecture.

\textbf{Formal Contracts and Drift.} \citet{bhardwaj2026abc} introduced Agent
Behavioral Contracts (ABC) applying Design-by-Contract to AI agents, with formal
drift bounds modeled as Ornstein-Uhlenbeck processes and (p, $\delta$,
k)-satisfaction guarantees. The drift detection in L4 shares motivation with ABC
and \citet{rath2026agent}'s multi-agent coordination metrics, but differs in
mechanism: L4 uses 12 mechanical features (mtime deltas, rule freshness, gate
failure rates) rather than behavioral divergence metrics, operating entirely
outside the Prose Barrier.

\textbf{Semantic Coordination.} \citet{acharya2026scf} proposed a 6-layer Semantic
Consensus Framework (SCF) for multi-agent systems, finding 79\% of failures stem
from coordination rather than capability. My dual-pool expert review
(Section~\ref{sec:dual-pool}) shares the concern about within-system consensus
reliability but focuses on single-agent persona independence (E1,
Section~\ref{sec:e1}) rather than cross-agent coordination.

\textbf{Governance Signals.} \citet{munirathinam2026inband} proposed the Recuse
Signal---a lightweight in-band governance standard (analogous to
\texttt{robots.txt}) for cooperative agent compliance. Access-time compliance
reached $\sim$82.5\%; mid-flight halts achieved 0/40. This supports a core
Prose Barrier implication: cooperative governance without enforcement is
insufficient for reliable constraint satisfaction.

\textbf{Format and Attention Routing.} \citet{pendereffect2026} provided
independent evidence that declarative vs. procedural format alters attention
routing. My L2/L3 experiments measure this at logprob resolution. The
non-monotonic constraint gradient extends this: the output-constraint-to-format
relationship is not linear.

\textbf{Self-Improvement vs Self-Verification.} \citet{wei2022cot} (chain-of-thought),
\citet{wang2022selfconsistency} (majority voting), \citet{madaan2023selfrefine}
(iterative self-feedback) show LLMs can self-improve. Why not self-verify? These
methods alter the generation distribution, not provide independent verification
channels. External verification (L1) and internal improvement (CoT) are
complementary, not substitutes.

\textbf{Critical Perspectives.} \citet{bender2020climbing}: language models
trained on form alone cannot acquire meaning. \citet{lipton2018troubling}:
overclaiming patterns in ML scholarship. These caution against interpreting
logprob shifts as ``deeper understanding'' and motivate explicit hedging.

\textbf{Independent Convergence.} During this work, Ren\'{e} Zander independently
built ``skillgate'' (L1 mechanical gates) without knowledge of this project---
convergent evolution supporting the architectural claim's generality. The
experimental agenda (P1-1, P1-2, E1) was shaped by DEV.to community threads:
Mike Czerwinski's receipt-of-action-vs-diligence taxonomy, Dipankar Sarkar's
gate-absent prediction, and Max Quimby's mechanizability boundary analysis.

\textbf{Gap.} No prior work combines (a) logprob-level measurement of constraint
penetration depth, (b) the non-monotonic constraint gradient (optimization
$\to$ suppression $\to$ rebound), (c) DPO-based behavioral internalization
measured via CIS, and (d) a unified five-layer architecture organized by position
relative to the Prose Barrier. TLOC proved formally what this work demonstrates
empirically; RVF/PolicyGuard/VIGIL enforce constraints externally but do not
measure internal constraint representation depth; ABC formalizes drift bounds
without logprob probes; SCF addresses multi-agent consensus without single-agent
barrier analysis. The Prose Barrier framework provides the missing integrative
lens.

\subsection{Practical Implications}

(1) Mechanically verifiable rules are fundamentally more enforceable. The
mechanizability scanner scores any candidate rule 0.0--1.0. (2) The $+0.38$
baseline compliance gain from having rules dwarfs format engineering effects---
invest in L1 first. (3) Light output constraints (C1) preserve the reasoning
channel; heavy suppression (C2) destroys it. The Prose Barrier of Measurement
is a structural constraint on LLM evaluation methodology.

\section{Limitations}
\label{sec:limitations}

\begin{enumerate}
\item \textbf{Single-rater, unblinded behavioral scoring (fatal).} $\kappa =
  -0.14$ on $n=8$. All behavioral numbers are exploratory. Independent blind
  scoring by $\geq$2 raters ($\kappa > 0.7$) is the critical precondition.

\item \textbf{Single-model logprob evidence.} Logprob V3 and constraint gradient
  conducted on DeepSeek V4 Pro only.

\item \textbf{Small samples.} Constraint gradient $n=12$/level. L0 Safety
  $n=5$ non-ceiling probes. Cross-model behavioral $n=12$, underpowered.

\item \textbf{Probe pre-screening selection bias.} \texttt{probe\_validator.py}
  selects probes where format effects already push tokens into measurable range.

\item \textbf{Constraint-length confound.} C2 vs C3 differ in length (21 vs 29
  Chinese characters).

\item \textbf{Post-hoc L1-visibility classification.} Single classifier, no
  second rater.

\item \textbf{Lack of pre-registration} (except P1-2).

\item \textbf{Single author, single device (Dell G15, RTX 3060 6GB), no
  supervision, no funding.}
\end{enumerate}

\section{Conclusion}

Agent configuration integrity is an architectural problem with a structural
root cause: the Prose Barrier. The five-layer architecture---permit (L0),
bypass (L1), detect (L2), encode (L3), predict (L4)---offers one systematic
response. Only L1 constitutes external verification; L2/L3 measure internal
representations.

The empirical evidence has clear limitations but points consistently: rules
work, format matters for internal processing but not behavior, and the output
constraint-to-format relationship is non-monotonic. The Prose Barrier of
Measurement has implications beyond this paper: binary-output experimental
designs may systematically underestimate format effects.

DPO experiments suggest behavioral preferences can be partially internalized
(CIS, $\sim$5--10\% NLL reduction), consistent with the Prose Barrier applying
at the weight level.

\section*{Community Acknowledgments}

This paper grew from 31 DEV.to articles written while trying to understand why
my agent kept ignoring its own rules. Thank you to the DEV.to community---
especially Mike Czerwinski (receipt-of-action vs receipt-of-diligence, paper-A
vs paper-B framing), Dipankar Sarkar (gate-absent prediction, decision-token
measurement), Max Quimby (mechanizability boundary), the anonymous reader who
spotted the $-10.0$ sentinel bug, René Zander (independent parallel discovery of
``skillgate''), and the seven commenters whose questions shaped every experiment,
including the ones that failed.

The remaining errors are mine alone. No supervisor reviewed this manuscript.
No funding supported this work.

\bibliographystyle{acl_natbib}
\bibliography{references}

\end{document}
